% \setlength{\columnsep}{12pt}
% espaço entre colunas %\setlength{\columnseprule}{0.2pt} % linha vertical entre colunas

%-------------------------------------------------------------------------
% Modelo Trabalho PPGE-UFPB/Estatística
% Normas da ABNT
% Provedor: Laboratório de Economia e Modelagem Aplicada -- LEMA
%-------------------------------------------------------------------------

\documentclass[12pt, openright,oneside, a4paper, english, brazil]{report}

% ----------------------------
% Pacotes básicos 
% ----------------------------
% Fonte e codificação (acentuação)
\usepackage{lmodern}

\renewcommand{\sfdefault}{ppl}
\usepackage[most]{tcolorbox}
\usepackage[T1]{fontenc}	
\usepackage[utf8]{inputenc}	
\usepackage[brazil]{babel}
\usepackage{comment} % para esconder a tabela Org do compilador
% Referências
%\usepackage[brazilian,hyperpageref]{}	
%\usepackage[alf]{abntex2cite}			
\usepackage[colon]{natbib}
\bibliographystyle{apalike}

%Texto aleatório
\usepackage{lipsum} 

% Tabelas e Figuras
\usepackage{ctable}             
\usepackage{multirow}
\usepackage{array}
\usepackage{float}              
\usepackage{longtable}          
\usepackage{xcolor,colortbl}    
\usepackage{booktabs}           
\usepackage{graphicx}			
\usepackage{subfig}             
\usepackage{caption}            
\usepackage{epstopdf}           

% Equações e símbolos
\usepackage{amsmath}            
\usepackage{nomencl} 
\usepackage{hyperref, cleveref}  

% Texto
\usepackage{indentfirst}		
\usepackage{color}				
\usepackage{microtype} 			
\usepackage{lastpage}			
% ----------------------------------------------------------
% Configurações do texto
% ----------------------------------------------------------
\usepackage{multicol}
\usepackage{array}
\newcolumntype{L}[1]{>{\raggedright\let\newline\\\arraybackslash\hspace{0pt}}m{#1}}
\newcolumntype{C}[1]{>{\centering\let\newline\\\arraybackslash\hspace{0pt}}m{#1}}
\newcolumntype{R}[1]{>{\raggedleft\let\newline\\\arraybackslash\hspace{0pt}}m{#1}}

\makeatletter
%\hypersetup{pdftitle={\@title}, pdfauthor={\@author}}
\makeatother

\def\preambulo{Trabalho elaborado em cumprimento às exigências da disciplina \disciplina, ministrada pela professor Dr. \professor.}
%--------------------------------------------------------------

\setlength{\parindent}{1.25cm}
\setlength{\parskip}{0.1cm} 

\usepackage{titlesec}
\usepackage{lipsum}

\makeindex




\title{Exercício Capitulo03}
\def\autor{Oscar Sérgio Monteiro}
\def\disciplina{SC334 - Avaliação de Desempenho Organizacional, turma b}
\def\professor{Antônio César Pitela}
\begin{document}
\makeatletter
\begin{titlepage}
    \begin{figure}[H]
    \vspace{-3cm}
    \centering
    \includegraphics[height=3cm,keepaspectratio]{config/coverpage}\hfill
    \par\vspace{1cm}
    \end{figure}
    \begin{figure}[H]
    \vspace{-2cm}
    \centering
    \includegraphics[height=3cm,keepaspectratio]{config/contabeis.png}
    \par\vspace{1cm}
    \end{figure}
     
    \vspace{2cm}
    
	\begin{center}
	    {\huge \bfseries \@title}
	    \begin{flushright}
		\begin{minipage}{8cm}
			\preambulo
		\end{minipage}
	\end{flushright}
	\vspace{3cm}
	    \normalsize
	    \autor
	\end{center}
	
\vfill
\centering
{\large Curitiba-PR, \the\year}

\end{titlepage}
\makeatother


\pdfbookmark[0]{\contentsname}{toc}
\tableofcontents
\newpage


\titleformat{\chapter}[display]{\normalfont\bfseries}{}{0pt}{\Huge  \thechapter. }
\newcommand{\doubletoprule}{%
  \specialrule{.1em}{2pt}{2pt}%
  \specialrule{.1em}{0pt}{0pt}%
}

\chapter{Caso WEG}
A WEG é frequentemente reconhecida pelo mercado pela combinação de:
\begin{itemize}
\item Crescimento sustentável;
\item Elevada rentabilidade;
\item Forte geração de caixa;
\item Controle do endividamento;
\item Internacionalização das operações
\end{itemize}
\section{Questões para Discussão}
\begin{itemize}
\item \textit{Quais indicadores financeiros seriam mais relevantes para analisara expansão internacional da empresa?}\\
O propósito dos indicadores financeiros é avaliar a saúde financeira de uma entidade, não vejo uma relação direta entre a saúde financeira e a expansão internacional. Importante seria verificar se os investimentos em expansão internacional, assim como outra qualquer estratégia de crescimento ou diversificação, estão impactando na saúde financeira da empresa, desta forma caberia avaliar se existe alguma relação entre a variação (melhora ou piora dos indicadores financeiros) face ao esforço de internacionalização. Entendo que os objetivos de internacionalização devem ser bem esclarecidos para que se definam indicadores de desempenho específicos para avaliar o atingimento dos objetivos. A medição dos indicadores financeiros, e por consequência a avaliação da saúde financeira, deve ser mantido e analisado sempre, havendo ou não novas iniciativas na entidade. Assim entendo que a expansão internacional deve ter seus indicadores de desempenho específicos. A analise das causas de eventual variação nos indicadores financeiros deve considera o textosto e o momento da empresa.
\item \textit{Como a geração de caixa influencia a capacidade de investimento?}\\
  A geração de caixa é um indicador financeiro dos mais relevantes, sem geração de caixa a entidade depende de capital de terceiros para realização de investimentos. A capacidade de investimento com capital próprio depende basicamente da geração de caixa, sem indicadores positivos de geração de caixa novos investimentos tem risco elevado. Todavia outros indicadores financeiros, como o nível de endividamento, também influenciam na capacidade de investimento.
\item \textit{Qual a relação entre rentabilidade e crescimento sustentável}\\
  Não existe uma relação direta entre rentabilidade de crescimento sustentável. Existem entidades com boa rentabilidade e crescimento insustentável. Uma maneira por exemplo, pode existir apenas para retirar a madeira para a formação de empreendimento agrícola. Uma vez extraída a madeira a empresa fecha. Mineradoras são outro exemplo. Encontramos mais situações de antagonismo entre crescimento sustentáveis e restabilidade em especial se a analise for de curto prazo. 
\end{itemize}
\section{Questões para debate}
\begin{itemize}
\item \textit{Uma empresa pode apresentar lucro e, simultaneamente,enfrentar dificuldades financeiras?}\\
  As finanças tem como elemento fundamental o valor no tempo, de forma que é possível uma entidade apresentas lucro em um exercício e simultaneamente ,mente ter alguma dificuldade financeira. Quando os insumos são paga a vista e as receitas a prazo. Esse descasamento de tempos pode gerar dificuldades financeiras.
\item \textit{O EBITDA é suficiente para avaliar o desempenho empresarial?}\\
A analise isolada de um ou poucos indicadores financeiro pode gerar uma avaliação com viés visto que cada indicador avalia um aspecto do desempenho. o EBITDA esta relacionada a margem de contribuição da empresa,não é suficiente para avaliar  desempenho\\
\item \textit{Em quais situações o ROE pode gerar interpretações equivocadas?}\\
Não apenas o ROE pode gerar interpretações equivocadas, em espacial se for analisado isoladamente. Em especial o ROE que compara a geração de lucro em relação ao capital próprio, pode ser enganoso quando a empresa que tenha um capital próprio pequena em relação aos seus ativos, pode apresentar um ROE positivo e não ter geração de caixa suficiente para o pagamento do juros do capital de terceiros.
\item {Como o endividamento afeta a rentabilidade dos acionistas?}\\
O endividamento geram despesas financeiras (não operacionais),esses custos que são pagos a terceiros antes dos pagamentos de dividendos. Empresas com muito capital de terceiros, que tenham um bom resultado operacional, são ótimas para os terceiros e não necessariamente para os sócios.
\item \textit{Quais são as limitações dos indicadores financeiros na economia digital?}\\
  Não vejo como os indicadores financeiros possam ser menos ilustrativos nas economias digitais, todos os indicadores devem ser analisados no contexto de cada negócio. Os negócios digitais tem suas peculiaridades, como por exemplo alta alavancagem por ter baixos custos fixos, essa característica deve ser considerada na analise dos indicadores. Um valor no índice de endividamento em uma empresa industrial pode ser considerado positivo, o mesmo valor em uma empresa digital poder ser muito negativo.
 \end{itemize}
\end{document}
